\chapter{Introducción}\label{introduccion}

En los últimos años, hemos visto como las tecnologías se disparaban. Desde el establecimiento de Internet y la World Wide Web todo lo informático ha crecido exponencialmente. De teléfonos móviles a teléfonos inteligentes hubo algo más de 20 años\footnote{Tomando de referencia el lanzamiento del DynaTAC\_8000X de Motorola en 1984 y el lanzamiento del primer IPhone de Apple en 2007.}, que es el mismo tiempo que ha pasado de este último a día de hoy.

Últimamente, no hemos dejado de oír de la computación cuántica. Cada poco tiempo se rompe una nueva barrera que previamente se pensaba que estaba a años de ser superada. Es por esto que la criptografía postcuántica es tan importante, en un mundo en el que nuestros datos personales vuelan de una aplicación a otra, proteger la infraestructura digital es no solo un deber sino una necesidad.

El National Institute of Standards and Technology o NIST es un organismo estadounidense perteneciente al ministerio de comercio que desde febrero de 2016, entre otros objetivos, busca establecer unos estándares de seguridad para que podamos mantener la misma seguridad de la que disfrutamos en la actualidad a pesar de este cambio de paradigma.

Debido a que toda migración a esta escala, es decir, a escala mundial, es costosa y lenta, una de las prioridades de la criptografía postcuántica es que todos estos algoritmos deban poder implementarse con los recursos que tenemos hoy en día y los protocolos de redes que utilizamos y, más seguramente, seguiremos utilizando incluso cuando existan los ordenadores cuánticos. Diseñar algoritmos y protocolos que soporten las capacidades de estas máquinas sin poder utilizar su potencia es un verdadero desafío pero no por ello menos importante.

Uno de los protocolos presentados a la convocatoria del NIST es el objeto de este trabajo, Super-singular Isogeny Key Exchange por Luca De Feo, David Jao y Jérôme Plût. SIKE es el único de los protocolos presentados que utiliza la composición de isogenias de curvas elípticas como base matemática de su funcionamiento.

SIKE fue parte de este proceso desde la primera ronda del NIST en diciembre de 2017 hasta la tercera que se cerró en julio de 2022. Sin embargo, poco después se descubriría un ataque efectivo sobre el protocolo que se formalizaría el año siguiente. Debido a esto, SIKE no se considera seguro y se retiraría del proceso de estandarizado \cite{status_report} .

A pesar de esta descalificación, hemos considerado que se trata de un protocolo interesante a estudiar debido a ser único en su base matemática y, precisamente, a su retirada del NIST. La retirada de este protocolo en concreto no implica que todo protocolo basado en las curva elípticas vaya a ser incompatible con la criptografía postcuántica, por lo que conocer las bases de la criptografía basada en isogenias no debería ser considerado impropio, sino necesario.

Este proyecto contiene una presentación de los conceptos matemáticos más relevantes, expuestos de manera sencilla y presentándolos progresivamente, para permitir la comprensión escalonada de estos. A continuación nos centraremos en el protocolo de SIKE, viendo el contexto en el que se desarrolla y su funcionamiento interno, así como su eventual descalificación del proceso del NIST.

El proyecto incluye también una interfaz gráfica, pensada de forma dual tanto para lectores de este documento como para informáticos interesados en la implementación en código, por lo que esta memoria también incluye toda la documentación técnica relacionada. Contamos entre otros el alcance, criterios de aceptación y requisitos, en los que definiremos los objetivos propios de la interfaz así como las condiciones bajo las cuales consideraremos que esta interfaz es válida. Contamos además de un listado completo de las tareas en las que determinamos el número total de horas invertidas en cada una y una especificación de la arquitectura, tecnologías y patrones de diseños que se han utilizado para el desarrollo. También incluimos un apartado de pruebas en el que hablaremos de la calidad del proyecto así como un apartado en el que desarrollaremos las incidencias y complicaciones que nos hemos encontrado y posibles desarrollos futuros que ampliarían la interfaz.
